\documentclass[conference]{IEEEtran}
\usepackage{iftex}
\ifPDFTeX
  \usepackage[utf8]{inputenc}
  \usepackage[T1]{fontenc}
\fi
\usepackage{amsmath,amssymb,bm,mathtools}
\usepackage{newtxtext,newtxmath}
\usepackage{booktabs,array,graphicx,siunitx,adjustbox}
\usepackage{placeins}
\usepackage{tikz,pgfplots}
\usetikzlibrary{arrows.meta,positioning,calc}
\pgfplotsset{compat=1.18}
\usepackage[final]{microtype}
\usepackage{cite}
\usepackage[breaklinks=true,hidelinks]{hyperref}
\usepackage[nameinlink,capitalise,noabbrev]{cleveref}
\interdisplaylinepenalty=2500
\emergencystretch=2em
\allowdisplaybreaks
\raggedbottom
\setcounter{topnumber}{5}
\setcounter{dbltopnumber}{5}
\renewcommand{\topfraction}{0.95}
\renewcommand{\dbltopfraction}{0.95}
\renewcommand{\textfraction}{0.05}
\renewcommand{\floatpagefraction}{0.8}
\renewcommand{\dblfloatpagefraction}{0.8}
\makeatletter
\setlength{\@fptop}{0pt}
\setlength{\@dblfptop}{0pt}
\makeatother
\providecommand{\mathdefault}[1]{#1}
\newcommand{\IncludePGF}[1]{\resizebox{\linewidth}{!}{\input{#1}}}
\newcommand{\vc}[1]{\bm{#1}}
\newcommand{\mt}[1]{\bm{#1}}
\newcommand{\R}{\mathbb R}
\newcommand{\C}{\mathbb C}
\newcommand{\E}{\mathbb E}
\newcommand{\tr}{^{\mathsf T}}
\newcommand{\herm}{^{\mathsf H}}
\newcommand{\sk}[1]{[#1]_\times}
\newcommand{\dd}{\mathrm d}
\DeclareMathOperator{\diag}{diag}
\DeclareMathOperator{\blkdiag}{blkdiag}
\DeclareMathOperator{\Var}{Var}
\DeclareMathOperator{\Rea}{Re}
\DeclareMathOperator{\Ima}{Im}
\DeclareMathOperator{\clip}{clip}
\pdfstringdefDisableCommands{\def\bm#1{#1}\def\vc#1{#1}\def\mt#1{#1}}

\hypersetup{pdftitle={Adaptive Shared-Mode CARMA Two-Frame Group},pdfauthor={Mikhail S. Grushinskiy}}
\title{Adaptive Shared-Mode CARMA Lie Group Filter for Marine Motion Estimation (draft)}
\author{\IEEEauthorblockN{Mikhail S. Grushinskiy}
\IEEEauthorblockA{Independent Researcher, Fair Lawn, NJ, USA\\
Bareboat Necessities GitHub Project\\mgrouch@users.github.com}}
\begin{document}\maketitle
\begin{abstract}
A vessel-mounted inertial unit measures specific force rather than wave
acceleration. Recovering surge, sway and heave therefore requires separating
orbital motion from gravity, attitude error and sensor bias. This article
presents the experimental 27-state shared-mode CARMA estimator on the
right-invariant two-frame Lie group implemented in ocean-IMU. A stable
CARMA(4,1) displacement prior supplies restoring dynamics without an
integral-displacement pseudo-measurement. Shared complex excitation carries
cross-axis phase; a minimal real realization retains six motion coordinates
per spatial axis. Three attitude errors and six body-bias errors complete
the 27-dimensional tangent covariance. We describe the spectral fit,
33-to-27 reduction, amplitude and frequency scaling, invariant prediction,
spatial process covariance, sensor correction and finite group reset.
An independent attitude proxy, spectral-moment period estimate, noise-aware
wave excitation, coupling confidence and staged magnetic startup determine
the causal adaptive schedule. Parameter changes transport the full posterior
through physical coordinates, retaining the attitude shear required by
unequal axis scales. Deterministic vessel-response replay supplies explicit
quality gates, with displacement, attitude and bias tracking. Existing paired
simulation evidence supplies vibration and sea-state-transition diagnostics.
The method estimates wave-relative motion, not absolute navigation; neither
universal hull accuracy nor global adaptive stability is assumed.
\end{abstract}
\begin{IEEEkeywords}
Marine inertial estimation, two-frame group, CARMA, heave, shared stochastic
modes, invariant error-state Kalman filter.
\end{IEEEkeywords}

\section{Estimation Objective and Information Paths}
\label{sec:objective}
The outputs are vessel attitude and displacement $\vc p$, velocity $\vc v$
and acceleration $\vc a$ associated with waves about a local equilibrium.
A local north--east--down (NED) frame is used. These sensors cannot identify
geographic position, steady translation, ocean current or a permanent
position offset merely by imposing a stochastic wave prior.

The estimator combines three kinds of structure. Offline spectral fitting
sets a fixed displacement-spectrum shape. A measurement-only frontend sets
its operating period and acceleration amplitude. A slower cross-spectral
branch estimates how axes share directional excitation. The two-frame group
specifies how attitude, world-frame motion and body-frame biases are corrected
together; it does not replace those stochastic or adaptation models.

The private attitude proxy runs from the first IMU sample, supplies startup
tilt, and levels acceleration for the period and energy estimators. Its
ongoing attitude is not fused as an extra sensor. The coupling branch has a
different information path: it uses measured specific force together with
the live estimated attitude and bias. Thus the period/energy branch avoids
main-filter feedback, while coupling estimation is deliberately state dependent.

The implementation reuses the same six-state motion realization, paired
process-noise discretizer, excitation mapping, confidence monitor and fusion
schedule as shared-mode quaternion estimation. A common TFG core provides
prediction, sensor updates and finite corrections. The following equations
specify the complete method without requiring the comparative study.

\section{Sensors, Tilt and Magnetic Information}
\label{sec:sensors}
Let $\mt R=\mt R_{bw}$ rotate body vectors into NED and set
$\mt R_{wb}=\mt R\tr$, $\vc g=g\vc e_z$. Calibrated sensor observations obey
\begin{align}
 \vc\omega_m&=\vc\omega+\vc b_g+\vc n_g,\label{eq:gyro}\\
 \vc f_m&=\mt R\tr(\vc a-\vc g)+\vc b_a+\vc k_a\Delta T+\vc n_a,\label{eq:acc}\\
 \vc m_m&=\mt R\tr\vc B_w+\vc n_m.\label{eq:mag}
\end{align}
The body-frame residual biases are $\vc b_g,\vc b_a$; $\vc k_a\Delta T$ is
an optional calibrated accelerometer temperature correction. Intrinsic sensor
calibration and physical alignment precede these equations. A configured
body lever arm $\vc r_b$ adds rotational acceleration
$\dot{\vc\omega}\times\vc r_b+
\vc\omega\times(\vc\omega\times\vc r_b)$ to specific force; the wrapper can
subtract it with the corresponding gyro-bias Jacobian. An unmodelled lever arm
makes the inferred translation describe the sensor location.

For illustration, a small body-frame world-to-body rotation error gives
\begin{equation}
 \delta\vc f=-\sk{\mt R_{wb}(\vc a-\vc g)}\delta\vc\theta
       +\mt R_{wb}\delta\vc a+\delta\vc b_a.
 \label{eq:ambiguity}
\end{equation}
A tilt error of order $\delta\theta$ can mimic horizontal acceleration of
order $g\delta\theta$ or trade against accelerometer bias. These are physical
information ambiguities, not problems solved merely by changing covariance
coordinates. Magnetic observations provide a direction independent of
instantaneous specific force. A usable horizontal field component supplies
heading information and helps distinguish tilt leakage from wave motion.
Field inclination, motion excitation, field-reference error and magnetic
disturbance still matter. During an outage, a simultaneous yaw rotation of
attitude and all world vectors leaves accelerometer predictions unchanged;
the wave prior is not an absolute yaw reference.
\section{Why a CARMA Displacement Prior?}
\label{sec:carma}
CARMA means continuous-time autoregressive moving-average. Its formal
scalar displacement equation is
\begin{equation}
 P(\mathcal D)p=g_d B(\mathcal D)\dot W,\qquad
 H_\kappa(s)=\frac{B_\kappa(s)}{P_\kappa(s)},
 \label{eq:carma}
\end{equation}
where $\mathcal D=\dd/\dd t$ and $\dot W$ is generalized white noise.
The state-space realization below, rather than differentiation of sampled
noise, defines its implementation. All roots of $P$ are strictly in the
left half-plane. For unit white-noise intensity the two-sided angular
frequency displacement spectrum is
\begin{equation}
 S_p(\omega)=g_d^2|H_\kappa(i\omega)|^2,\qquad
 \Var p=\frac{1}{2\pi}\int_{-\infty}^{\infty}S_p(\omega)\,\dd\omega.
 \label{eq:psd}
\end{equation}
Stable poles give decaying memory and a finite stationary displacement
variance. Complex-conjugate poles represent damped oscillatory response.
The moving-average numerator controls spectral zeros independently of the
restoring poles. A finite-dimensional state model then permits recursive
estimation at irregular sample intervals without storing a long motion
record. These are the relevant reasons to use CARMA here.

\subsection{Why CARMA(4,1)}
The implemented shape is
\begin{align}
 P_\kappa(s)&=L_{1,\kappa}(s)L_{2,\kappa}(s),\nonumber\\
 L_{j,\kappa}(s)&=s^2+2\zeta_j\kappa\omega_j s+(\kappa\omega_j)^2,
 \label{eq:poles}\\
 B_\kappa(s)&=s+\kappa\omega_0,\qquad \omega_0=0.
 \label{eq:zero}
\end{align}
Two damped pole pairs allow a narrow principal response and a broader
roll-off. The first-degree numerator gives relative degree $4-1=3$.
Consequently $S_p=O(|\omega|^{-6})$ at high frequency and the acceleration
spectrum $\omega^4 S_p$ is integrable: ideal acceleration has finite
variance. Acceleration itself is continuous in this realization, whereas
its derivative contains white driving noise. The fitted zero at the
origin gives $S_p=O(\omega^2)$ near DC and strongly discourages very slow
displacement. Four states are sufficient for this scalar spectral shape.

These features motivate the choice; they do not make CARMA(4,1) a unique
or exact law of ocean waves. In particular the ideal JONSWAP surface
spectrum has an $f^{-5}$ tail, so the fit is deliberately over a finite
band. Vessel response and sensor bandwidth also alter the actual motion
spectrum. The retained shape is a compact stochastic prior, not a hull
hydrodynamic model.

\subsection{Why the integral pseudo-measurement is absent}
An acceleration-driven chain with unrestrained velocity and position
integrators has poles at zero. One way to control its drift is to add
$S(t)=\int p(t)\,\dd t$ and repeatedly impose the soft constraint $S=0$.
In the present model, the restoring terms of the Hurwitz polynomial $P$
already make displacement stationary; there are no free position or
velocity integrators in the full state dynamics. The zero in
\cref{eq:zero} further suppresses displacement DC. An extra $S$ state and
an $S=0$ observation are therefore unnecessary for this regularization.

This does \emph{not} assert that the actual displacement is zero or that
an arbitrary initial displacement is observable. The prior chooses the
local wave equilibrium and assigns a distribution about it. It cannot
recover absolute position or eliminate tilt and bias ambiguities. A
covariance regularizer also cannot replace physical magnetic information.
The implementation allocates no $S_x,S_y,S_z$ states and executes no
integral-zero update; shared wrapper settings bearing that name have no
measurement effect in this configuration.

\section{The Offline Spectral Fit}
\label{sec:fit}
The prototype is fitted once to the JONSWAP spectral shape
\cite{jonswap}. Up to an amplitude constant, the target in frequency $f$ is
\begin{align}
 S_J(f)&\propto f^{-5}\exp\!\left[-\frac54(f_p/f)^4\right]
                  \gamma^{r(f)},\nonumber\\
 r(f)&=\exp\!\left[-\frac{(f-f_p)^2}{2\epsilon(f)^2f_p^2}\right],\label{eq:jonswap}\\
 \epsilon(f)&=\begin{cases}0.07,&f\le f_p,\\0.09,&f>f_p.\end{cases}\nonumber
\end{align}
The retained fit uses $f_p=0.29$~Hz, $\gamma=3.3$, and the uniform grid
$f_n=0.17,0.171,\ldots,1.30$~Hz. For trial shape parameters $\vartheta$,
\begin{align}
 d_n(\vartheta)&=10\log_{10}
       \frac{|H_1(i2\pi f_n;\vartheta)|^2}{S_J(f_n)},\nonumber\\
 J(\vartheta)&=\left[\frac1N\sum_{n=1}^N
                 (d_n-\bar d)^2\right]^{1/2}.
 \label{eq:fit}
\end{align}
Subtracting the mean $\bar d$ removes the arbitrary gain: this is a
\emph{shape} fit in decibels. It is not a likelihood fit to live IMU data,
nor a fit of displacement amplitude. The online acceleration statistic
sets the gain separately.

Bounded trust-region least squares uses 48 starts from a fixed random
seed. The constraints are
$f_1\in[0.15,0.60]$~Hz, $f_2\in[0.15,3.00]$~Hz,
$\zeta_1,\zeta_2\in[0.05,1]$, and
$f_0=\omega_0/(2\pi)\in[0,2]$~Hz. The optimum places $f_0$ at its lower
boundary. The stored implementation rounds that value to exactly zero
and uses the coefficients in \cref{tab:shape}.

\begin{table}[!t]
\caption{Prototype shape at frequency scale $\kappa=1$. Frequencies are
pole natural frequencies, not spectral peaks or zero-crossing frequencies.}
\label{tab:shape}\centering
\begin{tabular}{@{}lr@{}}\toprule
Parameter & Implemented value \\\midrule
$f_1=\omega_1/(2\pi)$ & 0.306945 Hz \\
$\zeta_1$ & 0.104191 \\
$f_2=\omega_2/(2\pi)$ & 0.815318 Hz \\
$\zeta_2$ & 0.478188 \\
$f_0=\omega_0/(2\pi)$ & 0 Hz \\\bottomrule
\end{tabular}
\end{table}

The unrounded fit has RMS shape error 1.856~dB on this grid and maximum
absolute residual 17.522~dB. The latter matters: a mean squared log fit
does not guarantee a uniformly small error, especially on the steep
low-frequency shoulder. Its displacement peak is about 0.307~Hz, while its
zero-crossing frequency is about 0.371~Hz. The operating-point adaptation
matches the latter, rather than identifying all three frequencies.

% Analytic curves only: no simulation CSV, external images, or shell escape.
\pgfmathdeclarefunction{carmaunit}{1}{%
 \pgfmathparse{(2*pi*#1)^2/((( (2*pi*0.306945)^2-(2*pi*#1)^2 )^2+(2*0.104191*(2*pi*0.306945)*(2*pi*#1))^2)*(((2*pi*0.815318)^2-(2*pi*#1)^2)^2+(2*0.478188*(2*pi*0.815318)*(2*pi*#1))^2))}}
\pgfmathdeclarefunction{jonswaplog}{1}{%
 \pgfmathparse{10/ln(10)*(-5*ln(#1)-1.25*(0.29/#1)^4+ln(3.3)*exp(-((#1-0.29)/(ifthenelse(#1<=0.29,0.07,0.09)*0.29))^2/2))}}
\begin{figure}[!t]\centering
\begin{tikzpicture}
\begin{axis}[width=\columnwidth,height=47mm,xmode=log,xmin=0.12,xmax=1.5,
 ymin=-65,ymax=3,xlabel={Frequency $f$ [Hz]},ylabel={Peak-normalized PSD [dB]},
 xtick={0.15,0.2,0.3,0.5,1,1.5},xticklabels={0.15,0.2,0.3,0.5,1,1.5},
 grid=major,tick label style={font=\footnotesize},label style={font=\small},
 legend style={font=\footnotesize,draw=none,at={(0.98,0.04)},anchor=south east}]
 \addplot[black,dashed,thick,domain=0.12:1.5,samples=180]
  {jonswaplog(x)-jonswaplog(0.29)};
 \addlegendentry{JONSWAP target}
 \addplot[blue!70!black,thick,domain=0.12:1.5,samples=180]
  {10*ln(carmaunit(x)/carmaunit(0.30736))/ln(10)};
 \addlegendentry{CARMA(4,1) shape}
\end{axis}
\end{tikzpicture}
\caption{The fitted rational spectral shape. Each curve is normalized to
its own peak for display; the fitting objective instead removes the mean
log-level difference over 0.17--1.30~Hz. The rational shape approximates a
finite band and does not reproduce the target's asymptotic tail.}
\label{fig:fit}
\end{figure}

\section{What Is Shared Between the Axes?}
\label{sec:shared}
Directional waves drive several components of vessel motion at once.
For example, horizontal orbital motion and vertical motion can be close
to quadrature: one reaches an extremum near the other's zero crossing.
Hull response adds its own phase. A real covariance between identical
scalar filters captures in-phase or opposite-phase dependence, but cannot
represent that quadrature relationship.

The method shares the \emph{temporal poles and stochastic excitation
structure}, not identical displacement trajectories. Each axis keeps its
own motion state and amplitude. The three driving processes have a
Hermitian covariance $\mt\Sigma=\mt\Sigma_R+i\mt\Sigma_I\succeq0$,
where $\mt\Sigma_R\tr=\mt\Sigma_R$ and
$\mt\Sigma_I\tr=-\mt\Sigma_I$. Real entries describe co-spectral
dependence; imaginary entries describe phase-shifted dependence. An
eigenfactorization of this matrix can be interpreted as a few common
random sources with different complex loadings on the axes. No such
eigenfactorization is required in the sample loop.

To construct the temporal mode, put one root from each pole pair into
\begin{align}
 P_+(s)&=(s-\lambda_{1+})(s-\lambda_{2+}),\nonumber\\
 \lambda_{j+}&=-\zeta_j\kappa\omega_j
         +i\kappa\omega_j\sqrt{1-\zeta_j^2}.
 \label{eq:pplus}
\end{align}
Write $P_+=N_r+iN_i$, with real-coefficient polynomials $N_r,N_i$.
Coefficient conjugation defines $P_-(s)=\overline{P_+(\bar s)}$; hence
$P=P_+P_-=N_r^2+N_i^2$. Choose the real, stable quadratic
$D(s)=s^2+d_1s+d_0$ so that
\begin{equation}
 |D(i\omega)|^2=|N_r(i\omega)|^2+|N_i(i\omega)|^2
 =\frac{|P_+(i\omega)|^2+|P_+(-i\omega)|^2}{2}.
 \label{eq:D}
\end{equation}
Its coefficients and realization are specified in \cref{app:colour}.
For circular complex driving increments $\dd\xi_i=\dd\xi_{r,i}+i\dd\xi_{i,i}$,
define the complex mode
\begin{equation}
 \zeta_i(s)=g_0\frac{B(s)}{P_+(s)D(s)}\xi_i(s),\qquad
 p_i=\Rea\zeta_i.
 \label{eq:complex}
\end{equation}
The increment normalization is
\begin{align}
 \E[\dd\vc\xi\,\dd\vc\xi\herm]&=2\mt\Sigma\,\dd t,\nonumber\\
 \E[\dd\vc\xi\,\dd\vc\xi\tr]&=0.
 \label{eq:circular}
\end{align}
Complex modes are an algebraic representation of real correlated motion;
they do not introduce complex-valued physical positions. Their connection
to multivariate CARMA and sums of complex Ornstein--Uhlenbeck modes is
discussed in \cite{fasen}. The particular factorization here is derived
explicitly rather than assumed from that general representation.

Let $H_c=B/(P_+D)$. The real displacement cross-spectrum is
\begin{equation}
 S_{ij}(\omega)=\frac{g_0^2}{2}
 \left[\Sigma_{ij}|H_c(i\omega)|^2+
       \overline{\Sigma_{ij}}|H_c(-i\omega)|^2\right].
 \label{eq:crossspectrum}
\end{equation}
Using \cref{eq:D}, its diagonal is exactly
$g_0^2\Sigma_{ii}|B/P|^2$. Its imaginary cross-spectrum is proportional to
$\Ima\Sigma_{ij}(|H_c(i\omega)|^2-|H_c(-i\omega)|^2)$ and is odd in
frequency, as required for real motion. Thus sharing preserves the chosen
scalar autospectra while allowing physically useful cross-axis phase.
It also restricts the possible frequency dependence: a single shared shape
and one band-estimated covariance are not a general model of several
independent crossing wave systems.

\section{From 33 States to 27}
\label{sec:reduction}
The complex transfer in \cref{eq:complex} has denominator degree four.
A direct complex state realization has four complex coordinates per
axis, or eight real coordinates. Three such axes plus three attitude
errors, three gyro biases and three accelerometer biases would give
$3\times8+3+3+3=33$ error states.

Only the real output and its derivatives enter the sensors. Rationalizing
the complex denominator yields the \emph{real}, two-input transfer
\begin{equation}
 \begin{aligned}
 p_i(s)&=g_0\frac{B(s)}{P(s)}\eta_i(s),\\
 \eta_i(s)&=\frac{N_r(s)\xi_{r,i}(s)+N_i(s)\xi_{i,i}(s)}{D(s)}.
 \end{aligned}
 \label{eq:realtransfer}
\end{equation}
Both inputs share the real denominator $PD$, of degree six. For the
implemented nondegenerate shape there is no common polynomial cancellation
between that denominator and both numerators. Its minimal real realization
therefore has six states, rather than eight.

The same conclusion follows from observability. In the eight-state
realification each root of $D$ occurs in both the complex mode and its
conjugate. One real displacement output cannot distinguish both directions
in each repeated eigenspace. The corresponding two-dimensional invariant
subspace has identically zero displacement, velocity and acceleration
readout. Removing it eliminates two states per axis, hence six overall.
Rotating the acceleration readout with attitude does not reveal a direction
whose three acceleration outputs are all zero. Magnetic observations act
on attitude, not on those hidden motion directions.

This is a minimal realization of the motion law at fixed parameters. With
compatible projected initial distributions it preserves the complete
three-axis observable Gaussian law, including quadrature spectra and the
corresponding linearized Kalman updates. It is not a universal claim that
arbitrary time-varying coordinates in an unreduced model give the same
adaptive trajectories. Parameter-change transport must be specified; the
actual six-state cascade does so in \cref{sec:commits}.

\subsection{The six-state real cascade}
Set $\bar\omega_j=\kappa\omega_j$,
$u_j=2\zeta_j\bar\omega_j$ and $d=-u_2$ for the implemented zero
$\omega_0=0$. The four scalar spectral states are
$\vc z_i=[p_i,v_i,a_i,\ell_i]\tr$, with
\begin{equation}
 \mt A_4=\begin{bmatrix}
 0&1&0&0\\0&0&1&0\\
 d\bar\omega_1^2&du_1-\bar\omega_1^2&d-u_1&-\bar\omega_2^2\\
 \bar\omega_1^2&u_1&1&0
 \end{bmatrix}.
 \label{eq:A4}
\end{equation}
Its input acts on acceleration through $\vc e_a=[0,0,1,0]\tr$.
The coordinate $\ell_i$ is spectral memory: it is a latent state of the
forcing oscillator in a factored realization. Its equation is
$\dot\ell_i=\bar\omega_1^2p_i+u_1v_i+a_i$.
In particular $\ell_i$ is not an extra $S_i=\int p_i\,\dd t$ state on
which zero is observed.

The two additional states $\vc w_i\in\R^2$ realize the colouring row
\begin{equation}
 \mt V(s)=\frac{[N_r(s)\quad N_i(s)]}{D(s)}
       =[1\quad0]+\mt C_w(s\mt I-\mt A_w)^{-1}\mt B_w.
 \label{eq:V}
\end{equation}
Equation~\eqref{eq:D} implies
$\mt V(i\omega)\mt V(i\omega)\herm=1$. This property, called
\emph{co-inner} or lossless, means that independent unit white inputs give
a unit white output in the stationary Gaussian law. The two states supply
phase structure without changing the single-axis autospectrum.

Use input-normal coordinates, whose unit-input controllability Gramian is
$\mt I_2$:
\begin{equation}
 \mt A_w+\mt A_w\tr=-\mt B_w\mt B_w\tr,\qquad
 \mt C_w=-[1\quad0]\mt B_w\tr.
 \label{eq:inputnormal}
\end{equation}
The final axis state and dynamics are
\begin{align}
 \vc x_i&=[p_i,v_i,a_i,\ell_i,w_{i1},w_{i2}]\tr,\nonumber\\
 \mt A_6&=\begin{bmatrix}\mt A_4&\vc e_a\mt C_w\\0&\mt A_w\end{bmatrix},&
 \mt G_6&=\begin{bmatrix}\vc e_a[1\quad0]\\\mt B_w\end{bmatrix},\label{eq:A6}\\
 \dd\vc x_i&=\mt A_6\vc x_i\,\dd t+
                 g_0\mt G_6\begin{bmatrix}\dd\xi_{r,i}\\\dd\xi_{i,i}\end{bmatrix}.
 \nonumber
\end{align}
There is direct noise only in acceleration and colouring memory;
$\dot p_i=v_i$ and $\dot v_i=a_i$ remain exact kinematic relations.


\begin{table}[!t]\centering
\caption{Invariant tangent allocation. Nominal attitude is a unit quaternion
(or rotation matrix), with three attitude errors. World columns are stored
column-major; the covariance has 27 coordinates.}\label{tab:states}
\begin{tabular}{@{}llr@{}}\toprule
Block&Coordinates&Count\\\midrule
Attitude&$\vc\phi$&3\\
Gyro bias&$\vc\beta_g$&3\\
Displacement, velocity, acceleration&$\vc\rho_p,\vc\rho_v,\vc\rho_a$&9\\
Spectral and colouring memory&$\vc\rho_\ell,\vc\rho_{w_1},\vc\rho_{w_2}$&9\\
Accelerometer bias&$\vc\beta_a$&3\\\midrule
Total&&27\\\bottomrule
\end{tabular}\end{table}


\section{Scaling in Time and Amplitude}
\label{sec:scaling}
Let the unit-gain, unit-scale moments of the four-state shape be
$M_p=\Var_1p$, $M_v=\Var_1v$ and $M_a=\Var_1a$. They are obtained once
in double precision from
\begin{equation}
 \mt A_4(1)\mt P_1+\mt P_1\mt A_4(1)\tr+
                  \vc e_a\vc e_a\tr=0.
 \label{eq:basecov}
\end{equation}
Scaling every pole and zero by $\kappa$ gives
\begin{equation}
 H_\kappa(s)=\kappa^{-3}H_1(s/\kappa).
 \label{eq:similarity}
\end{equation}
Changing integration variable to $\omega/\kappa$ in \cref{eq:psd} yields
\begin{equation}
 \Var p=\frac{g_d^2M_p}{\kappa^5},\qquad
 \Var v=\frac{g_d^2M_v}{\kappa^3},\qquad
 \Var a=\frac{g_d^2M_a}{\kappa}.
 \label{eq:momentscale}
\end{equation}
To request acceleration standard deviation $\sigma_i$, use
\begin{equation}
 g_0^2=\frac{\kappa}{M_a},\qquad \Sigma_{ii}=\sigma_i^2.
 \label{eq:gain}
\end{equation}
Thus $g_d^2=g_0^2\sigma_i^2$ for each marginal, and
\begin{equation}
 \sigma_{p_i}=\frac{\sigma_i}{\kappa^2}\sqrt{M_p/M_a},\qquad
 \sigma_{v_i}=\frac{\sigma_i}{\kappa}\sqrt{M_v/M_a}.
 \label{eq:motionstd}
\end{equation}
At a fixed acceleration scale, doubling the wave frequency quarters the
displacement scale. Doubling $\sigma_i$ doubles displacement and velocity
standard deviations and quadruples all marginal variances.

The prototype displacement zero-crossing frequency is
\begin{equation}
 f_{z,1}=\frac{1}{2\pi}\sqrt{M_v/M_p},\qquad
 \kappa=\frac{f_{\rm op}}{f_{z,1}}.
 \label{eq:kappa}
\end{equation}
This uses the same spectral-moment quantity estimated online. It does not
equate measured period to the reciprocal of a pole natural frequency.
All three axes use the same $\kappa$ and shape. The uncoupled scheduled amplitudes are
isotropic, $\sigma_x=\sigma_y=\sigma_z$ (the coherence cap can enlarge
horizontal excitation); there is no assumption that their
sample paths are identical.

The colouring realization also scales exactly:
\begin{equation}
 \mt A_w(\kappa)=\kappa\mt A_{w1},\quad
 \mt B_w(\kappa)=\sqrt\kappa\mt B_{w1},\quad
 \mt C_w(\kappa)=\sqrt\kappa\mt C_{w1}.
 \label{eq:colourscale}
\end{equation}
Only the unit-scale two-state factorization is needed. Its input-normal
Gramian remains $\mt I_2$ at every frequency scale.


\section{Two-Frame State and Right-Invariant Error}
\label{sec:geometry}
The nominal state is $X=(R,W,B)$ with
\begin{equation}
 W=[\vc p,\vc v,\vc a,\vc\ell,\vc w_1,\vc w_2]\in\R^{3\times6},
 \qquad B=[\vc b_g,\vc b_a]\in\R^{3\times2}.
\end{equation}
The two-frame product and inverse are
\begin{align}
 X_1X_2&=(R_1R_2,\ W_1+R_1W_2,\ B_2+R_2\tr B_1),\\
 X^{-1}&=(R\tr,-R\tr W,-RB).
\end{align}
World columns transform on the left and body biases on the right. This is
$\operatorname{TwoFrameGroup}(3,6,2)$, following the navigation geometry in
\cite{barrau}. It has $3+3(6)+3(2)=27$ tangent dimensions. Its error is
\begin{align}
 \eta&=\hat X X^{-1}=\operatorname{Exp}_G(\xi),\\
 \xi&=(\vc\phi,\vc\beta_g,\vc\rho_p,\vc\rho_v,\vc\rho_a,\nonumber\\
 &\qquad\vc\rho_\ell,\vc\rho_{w_1},\vc\rho_{w_2},\vc\beta_a).
\end{align}
To first order, $\delta R=\hat R R\tr$,
$\vc\rho_i=\hat{\vc w}_i-\delta R\vc w_i$ and
$\vc\beta=R(\hat{\vc b}-\vc b)$. Thus the bias-error coordinates are
world-referred even though the nominal biases are body-referred. Group
exponential and logarithm use the SO(3) Jacobians on vector slots; they are
not componentwise addition and subtraction.

For physical additive errors with true minus estimated means,
$e=(\delta\theta_b,\delta b_g,\delta w_0,\ldots,\delta w_5,\delta b_a)$,
the local coordinate map $\xi=T_{\hat X}e$ is
\begin{align}
 \vc\phi&=\hat R\delta\vc\theta_b,&\vc\beta&=-\hat R\delta\vc b,\\
 \vc\rho_i&=-\delta\vc w_i+\sk{\hat{\vc w}_i}\hat R\delta\vc\theta_b.
 \label{eq:physicalmap}
\end{align}
This map specifies how a physical uncertainty seed, a bias hold or a motion
coordinate change is expressed in invariant covariance. A numerically equal
matrix in the two coordinate systems does not represent equal uncertainty.

\section{Invariant Prediction and Discrete Process Covariance}
\label{sec:kalman}
Let $\Pi$ be the full invariant tangent covariance. The deterministic
nominal dynamics are
\begin{align}
 \dot R&=R\sk{\vc\omega_m-\hat{\vc b}_g},&\dot W&=WA_6\tr,\\
 \dot{\vc b}_g&=0,&\dot{\vc b}_a&=-\vc b_a/\tau_b.
\end{align}
When the common startup schedule holds accelerometer bias, its last equation
is replaced by zero and its bias-driving covariance is held as well. The
linearized invariant equations, with $\vc\omega_W=\hat R(\vc\omega_m-\hat b_g)$,
are
\begin{align}
 \dot{\vc\phi}&=-\vc\beta_g+\hat R\vc n_g,\\
 \dot{\vc\beta}_g&=\sk{\vc\omega_W}\vc\beta_g-\hat R\vc n_{bg},\\
 \dot{\vc\rho}_i&=\sum_j A_{6,ij}\vc\rho_j-\sk{\hat{\vc w}_i}\vc\beta_g\nonumber\\
 &\qquad+\sk{\hat{\vc w}_i}\hat R\vc n_g-\vc n_{c,i},\\
 \dot{\vc\beta}_a&=(\sk{\vc\omega_W}-\tau_b^{-1}I)\vc\beta_a-\hat R\vc n_{ba}.
 \label{eq:invariantcontinuous}
\end{align}
The common dynamics on all three axes commute with rotation, cancelling a
deterministic attitude-to-motion term. The remaining bias-to-motion term
depends on estimated world columns. Accordingly this is a right-invariant
error-state EKF, not an assertion of exact autonomous log-linear error dynamics.

For interval $h$, let $F_c=e^{A_6h}$, $R_1=R_0\operatorname{Exp}([\omega h]_\times)$,
$E=R_1R_0\tr$, $F_b=-R_0hJ_l(\omega h)R_0\tr$, and $W_1=W_0F_c\tr$.
The nonzero discrete error blocks, in addition to the attitude identity, are
\begin{align}
 \Phi_{\phi,b_g}&=F_b,&\Phi_{b_g,b_g}&=E,\\
 \Phi_{\rho_i,\rho_j}&=F_{c,ij}I_3,&
 \Phi_{\rho_i,b_g}&=\sk{w_{1,i}}F_b,\\
 \Phi_{b_a,b_a}&=\phi_b E,&\phi_b&=e^{-h/\tau_b}\ \text{(active)},
\end{align}
with $\phi_b=1$ while held. Prediction is $\Pi^-=\Phi\Pi^+\Phi\tr+Q$,
retaining every cross-covariance.

The spatial excitation is $\Sigma=S_R+iS_I$, $S_R=S_R\tr$, $S_I=-S_I\tr$.
For $J_2=\left[\begin{smallmatrix}0&-1\\1&0\end{smallmatrix}\right]$,
\begin{align}
 D_R&=g_0^2G_6G_6\tr,&D_I&=g_0^2G_6J_2G_6\tr,\\
 Q_R(h)&=\int_0^h e^{A_6t}D_Re^{A_6\tr t}\,\dd t,\\
 Q_I(h)&=\int_0^h e^{A_6t}D_Ie^{A_6\tr t}\,\dd t.
\end{align}
The motion covariance between world columns is the $3\times3$ block
\begin{equation}
 Q^{\rm motion}_{ij}=Q_R(i,j)S_R+Q_I(i,j)S_I.
 \label{eq:Qpair}
\end{equation}
$Q_R$ is symmetric and $Q_I$ antisymmetric. Opposite spatial blocks transpose,
so the complete real covariance is symmetric and positive semidefinite.
The imaginary term carries quadrature coupling; deleting it changes the
cross-axis process even if marginal spectra look unchanged.

The paired integrals are evaluated by the shared scaled Taylor discretizer
and repeated doubling, a continuous-to-discrete covariance construction
\cite{vanloan}. At small $h$, series are accumulated to working precision;
larger intervals use
\begin{equation}
 Q(2h)=Q(h)+F_c(h)Q(h)F_c(h)\tr,\quad F_c(2h)=F_c(h)^2.
\end{equation}
A cached reference and local step-size series handle ordinary timestamp
jitter. The continuous-time expressions are exact; their implementation uses
finite series and floating-point arithmetic.

Motion uncertainty also receives gyro and bias-driving noise through the
estimated columns. The identity
$F_c(h-t)W(t)\tr=W(h)\tr$ collapses the gyro impulse response. In particular,
\begin{align}
 Q_{\rho_i,\phi}&=\sk{w_{1,i}}Q_{\phi\phi},\\
 Q_{\rho_i,b_g}&=\sk{w_{1,i}}Q_{\phi b_g},\\
 Q^{\rm gyro}_{\rho_i,\rho_j}&=
       \sk{w_{1,i}}Q_{\phi\phi}\sk{w_{1,j}}\tr.
\end{align}
The small remaining rate/bias integrals use five-point Gauss--Legendre
quadrature. Structured two-pass covariance propagation reuses the six-state
transition; it does not replace these correlations with an isolated motion
covariance or require a dense 27-state matrix exponential.

\section{Sensor Updates and Finite Group Reset}
\label{sec:updates}
The world-referred residuals are
\begin{align}
 r_a&=\hat R(f_m-\hat b_a-k_a\Delta T)-(\hat a-g),\\
 r_m&=\hat Rm_m-B_w.
\end{align}
In the tangent ordering above their Jacobians, absent a lever arm, are
\begin{align}
 H_a&=[\sk g\ \ 0\ \ 0\ \ 0\ \ {-I}\ \ 0\ \ 0\ \ 0\ \ {-I}],\\
 H_m&=[-\sk{B_w}\ \ 0\ \cdots\ \ 0].
\end{align}
Sensor covariance is rotated into the residual frame,
$R_y=\hat R R_{y,b}\hat R\tr$. A configured lever arm adds
$H_{a,\beta_g}=-\hat R J_L\hat R\tr$ for the derivative of rotational
acceleration with respect to body gyro bias. The state-independent Jacobian
statement therefore has this explicit optional exception.

For either sensor,
\begin{align}
 \Omega&=H\Pi^-H\tr+R_y,&K&=\Pi^-H\tr\Omega^{-1},\\
 d&=Kr,&\hat X^+&=\operatorname{Exp}_G(-d)\hat X^-.
\end{align}
\begin{align}
 \Pi^J&=(I-KH)\Pi^-(I-KH)\tr+KR_yK\tr,\\
 \Pi^+&=J_r(d)\Pi^J J_r(d)\tr,\label{eq:groupreset}\\
 J_r(d)&=\left.\frac{\partial}{\partial\epsilon}
 \operatorname{Log}_G\bigl(\operatorname{Exp}_G(-d)
                \operatorname{Exp}_G(d+\epsilon)\bigr)\right|_0.
\end{align}
Small systems are factorized rather than explicitly inverted. Group injection
corrects all world columns and both body biases along with attitude. The full
right Jacobian transports all corresponding covariance blocks. Replacing it
by an attitude-only approximation or identity is a different estimator,
retained only as a research ablation.

The common wrapper's frozen-bias semantics and physical bias bounds remain
in force. Its accelerometer NIS, $r_a\tr\Omega^{-1}r_a$, feeds coupling
confidence. This monitor is not a per-packet magnetic outlier gate. Finite
nonzero magnetic measurements with a successful innovation factorization are
accepted by the inherited core; this article does not claim additional
chi-square disturbance rejection or guaranteed informative magnetic service.
\section{The Independent Proxy and the Wave-Period Estimator}
\label{sec:proxy}
\subsection{Why use a private attitude proxy}
The frontend runs a private gyro--accelerometer complementary observer
of Mahony type \cite{mahony}. It runs from the first sample and produces
up-positive levelled vertical acceleration
\begin{equation}
 a_u=-\vc e_z\tr\mt R_{bw,\mathrm{proxy}}\vc f_m-g.
 \label{eq:verticalproxy}
\end{equation}
Its accelerometer correction corner lies below the wave band; gyro
propagation carries attitude through orbital acceleration. The configured
gain parameters are \texttt{two\_kp}$=0.2$ and
\texttt{two\_ki}$=0.02$. The proportional corner is approximately
0.016~Hz. A fast gravity alignment would follow wave acceleration into a
false tilt. The proxy's yaw is unobservable and may drift; it is removed
when that observer is used as a magnetic accumulation frame.

This branch supplies startup tilt and the default period and energy
measurements. It does not use the main filter's attitude, estimated
acceleration, displacement or bias. Otherwise a main-filter error could
change the apparent sea state, retune its own prior, and reinforce itself.
After startup the proxy continues to supply the tuner; it is not fused as
an additional attitude sensor. The three-axis coupling estimator is a
separate branch and \emph{does} use the main attitude and bias.

\subsection{Period from spectral moments, not acceleration peaks}
For stationary differentiable displacement with moments
$m_0=\Var p$ and $m_2=\Var\dot p$, the spectral zero-crossing period is
\begin{equation}
 T_z=2\pi\sqrt{m_0/m_2}.
 \label{eq:Tz}
\end{equation}
This is generally neither the spectral peak period nor the period of the
acceleration-band carrier. Acceleration weights a displacement spectrum
by $\omega^4$, emphasizing faster components.

The frontend estimates \cref{eq:Tz} with two high-pass stages followed
by two leaky integrations of $a_u$. With leak rate
$\lambda=2\pi(0.02)$~s$^{-1}$, their continuous-time transfer analogues are
\begin{equation}
 \tilde v(s)=\frac{s^2}{(s+\lambda)^3}a_u(s),\qquad
 \tilde p(s)=\frac{s^2}{(s+\lambda)^4}a_u(s).
 \label{eq:periodfilters}
\end{equation}
The actual recursions use timestamp-dependent first-order exponential
coefficients. Bias-corrected exponentially weighted first and second
moments give variances $V_v,V_p$, followed by
\begin{equation}
 T_{z,\mathrm{raw}}=\frac{2\pi}{\sqrt{V_v/V_p-\lambda^2}}.
 \label{eq:Traw}
\end{equation}
The leak subtraction is exact for a narrowband sinusoid in the
continuous-time analogue. For broadband motion it gives a filter-weighted
moment ratio, not an exact unfiltered $T_z$.

The moment horizon is
$\clip(4\hat T_z,20,180)$ seconds. A single canonical log-period state
smooths $\log T_{z,\mathrm{raw}}$ with horizon $0.05\hat T_z$ subject to
implementation guards. Period and frequency are the exponentials of that
state with opposite signs, so they remain exact reciprocals; there is no
second frequency smoother in this chain.

Statistics start after $3/\lambda\simeq23.9$~s to avoid the leaky
integrator startup transient. A usable period requires elapsed time at
least $4/\lambda\simeq31.8$~s and at least one current period of moment
history. This takeover is latched. The stricter settled diagnostic also
requires $6/\lambda\simeq47.7$~s and an accumulated moment weight above
one half. Until usable, the tuner uses a fixed 0.2~Hz prior. These clocks
and energy confidence, rather than the mere existence of a numerical
ratio, determine whether the period can guide adaptation.

\section{Noise-Aware Amplitude and Frequency Adaptation}
\label{sec:adapt}
\subsection{Band energy above the sensor-noise prediction}
The vertical channel is passed through a period-scaled band with
$f_l=0.5 f_t$ and $f_h=4f_t$, limited to 0.01--6~Hz. Its continuous-time
analogue is
\begin{equation}
 H_B(s)=\frac{s b}{(s+a)(s+b)},\qquad a=2\pi f_l,\quad b=2\pi f_h.
 \label{eq:band}
\end{equation}
The band uses the preceding frequency value so the current observation
does not choose its own observation bandwidth. The variance horizon is
approximately four current periods, with bounds and sample-time guards.
Matched exponentially weighted first and second moments give
$V=\E_w[y^2]-(\E_w[y])^2$ after normalization by their accumulated weight.

Configured sensor noise consists of white per-sample standard deviation
$s_n$, bias random-walk intensity $q_b$, and relative uncertainty $u$ in
the predicted noise variance. The no-wave variance is estimated as
\begin{equation}
 V_0=s_n^2G_B+
 \frac{q_b b}{2a(a+b)}-\frac{q_b}{2T_Va^2},
 \label{eq:V0}
\end{equation}
clipped nonnegative. $G_B$ is the band's tracked discrete white-noise
variance gain, and $T_V$ is the variance horizon. The last two terms
approximate the random-walk contribution and removal of its slow mean.
The effective degrees of freedom $\nu$ account for band correlation,
finite horizon and initial accumulated weight; their implemented expression
is given in \cref{app:evidence}.

Define excess variance $E=V-V_0$ and a dimensionless detection statistic
\begin{equation}
 z_E=\frac{E}{V_0\sqrt{2/\nu+u^2}},\qquad
 p_E=\mathcal S\!\left(\frac{z_E-2}{4-2}\right),
 \label{eq:evidence}
\end{equation}
where $\mathcal S(x)=0$ for $x\le0$, $3x^2-2x^3$ for $0<x<1$, and
1 for $x\ge1$. The default evidence-mode excitation target is
\begin{equation}
 \sigma_{\rm target}=c_\sigma\sqrt{p_E\max(E,0)}.
 \label{eq:sigmatarget}
\end{equation}
Although the API calls $p_E$ a probability, it is a configured smooth
evidence weight, not a calibrated Bayesian posterior probability.
No physical wave variance is assigned merely because the sensor has a
noise floor. The numerical acceleration floor, $10^{-5}$~m/s$^2$, exists
only to keep model configuration and floating-point covariance valid.

The evidence branch starts its own band copy at the first acceleration
sample rather than at zero. That avoids interpreting a DC turn-on step as
a minute of wave energy. Its white-noise and bias settings must describe
the installed device. The defaults, $s_n=0.0148$~m/s$^2$ at 200~Hz,
$\sqrt{q_b}=5\times10^{-4}$~m/s$^2/\sqrt{\mathrm s}$ and $u=0.5$, are
simulation settings, not a universal sensor calibration.

\subsection{Rejecting a period unsupported by acceleration energy}
The doubly integrated period proxy emphasizes low-frequency errors.
For small waves, a slow bias or levelling transient can suggest a period
longer than that of the actual acceleration band. An extra low-pass at
$f_h$ forms $y_2$, and a finite-difference derivative forms $\dot y_2$.
After subtracting their modeled no-wave variances, define
\begin{equation}
 f_c=\frac1{2\pi}\sqrt{\frac{V_{\dot y_2}-V_{\dot y_2,0}}
                                  {V_{y_2}-V_{y_2,0}}},\qquad
 r_f=\frac{f_c}{\sqrt{f_lf_h}}.
 \label{eq:fc}
\end{equation}
Only positive excess moments yield a valid ratio. The period confidence is
\begin{equation}
 c_f=\mathcal S\!\left(\frac{z_E-8}{16-8}\right)
         \mathcal S\!\left(\frac{2.5-r_f}{2.5-2.0}\right).
 \label{eq:cf}
\end{equation}
A band centred far below the waves puts its excess energy near the upper
edge, increasing $r_f$ and reducing this confidence. This tests the
frequency support of the period; it is not a new full-spectrum fit.

\subsection{Smoothing and staged activation}
The common scheduling coordinate is $\tau=c_\tau T_z/2$, with default
$c_\tau=1$. Here $\tau$ parameterizes the operating period; it is not an
additional OU acceleration state in the motion model. Frequency and
$\tau$ are clamped to their configured operating envelope. With
$\alpha=1-e^{-h/T_A}$, default $T_A=0.4(T_z/2)=0.2T_z$ subject to guards,
the per-sample candidates obey
\begin{align}
 \tau^+&=\tau+\alpha w_f(\tau_{\rm target}-\tau),\nonumber\\
 \sigma^+&=\sigma+\alpha(\sigma_{\rm target}-\sigma),\label{eq:ema}\\
 w_f&=p_Ec_f\quad\text{once the measured period is usable}.\nonumber
\end{align}
Before takeover, $w_f=1$ follows the fixed period prior. Later, no
supporting wave evidence means $w_f=0$: frequency is held while amplitude
can decay toward the numerical floor. The internal coefficient
$c_\sigma=0.9$ is divided out when assigning acceleration amplitude,
\begin{equation}
 f_{\rm op}=\frac{c_\tau}{2\tau},\qquad
 \sigma_i=c_{\rm gain}\max(10^{-5},\sigma/c_\sigma),
 \label{eq:operatingpoint}
\end{equation}
with $c_{\rm gain}=1$ and equal axes by default.

Every valid IMU sample contributes to the moments and smoothing.
Activation is throttled separately: about every 0.1~s a candidate is
staged, then committed at the start of the next IMU sample. Thus a sample
is predicted and corrected with previously active model parameters.
Changing process coefficients does not reset the posterior to a fresh
stationary covariance on each commit.

\section{Estimating Shared Excitation}
\label{sec:crossestimate}
The coupling branch receives measured world-frame acceleration,
\begin{equation}
 \vc a_{\rm meas}=\hat{\mt R}_{bw}(\vc f_m-\hat{\vc b}_a)+\vc g,
 \label{eq:ameas}
\end{equation}
with temperature and lever-arm corrections when configured. It does not
use estimated CARMA acceleration as its input, but it does depend on the
estimated attitude and bias. Consequently it is not statistically
independent of the main filter.

Complex demodulation at $\dot\phi=2\pi f_{\rm op}$ gives
\begin{equation}
 A_i=\mathcal L\{a_{{\rm meas},i}e^{-i\phi}\},\qquad
 \hat{\mt S}=\E_w[\vc A\vc A\herm].
 \label{eq:demod}
\end{equation}
$\mathcal L$ is two cascaded first-order low-passes with corner
$0.15f_{\rm op}$. The outer-product average has horizon 240~s and is
normalized by accumulated weight. It is Hermitian positive semidefinite
by construction. A practical effective sample count is
$N_{\rm eff}=\max(1,2\min(t,240)\,0.15f_{\rm op})$.
Coupling is not eligible before 120~s of this branch's data.

At an individual frequency, $a_i=-\omega^2p_i$ on all axes, so
$S_{a_i a_z}/S_{a_z a_z}=S_{p_i p_z}/S_{p_zp_z}$.
Within the finite demodulation band this is an approximation when the
response ratios vary with frequency. The estimate describes local
cross-spectral regressions, not a measured full-band hydrodynamic transfer
function.

Partition the band matrix into horizontal and vertical blocks. Its
multiple coherence is estimated as
\begin{equation}
 \gamma^2=\clip\!\left(
 \frac{\hat{\vc S}_{hz}\herm\hat{\mt S}_{hh}^{-1}
                         \hat{\vc S}_{hz}}{\hat S_{zz}},0,1\right).
 \label{eq:coherence}
\end{equation}
If the horizontal block is nearly singular, the implementation uses its
trace-normalized cross-power surrogate instead. Statistical shrinkage is
$s_c=\max(0,1-1/(N_{\rm eff}\gamma^2))$, or zero if coherence is zero.
With coupling confidence $c\in[0,1]$ and user gain $g_c$, set
\begin{equation}
 \vc b=cg_cs_c\frac{\sigma_z^2}{\hat S_{zz}}\hat{\vc S}_{hz},\qquad
 \mt\Sigma=\begin{bmatrix}\mt H&\vc b\\\vc b\herm&\sigma_z^2\end{bmatrix},
 \label{eq:excitation}
\end{equation}
where the default $\mt H=\sigma_h^2\mt I_2$ retains the scheduled
horizontal amplitudes. No measured horizontal autospectrum is promoted
directly to wave energy: horizontal tilt leakage would contaminate it.
The real and imaginary horizontal--vertical terms are retained.
If $q=\vc b\herm\mt H^{-1}\vc b>\rho\sigma_z^2$ ($\rho=0.8$), enlarge
$\mt H$ along the coupling only,
$\mt H\leftarrow\mt H+\frac{\lambda-1}{q}\vc b\vc b\herm$ with
$\lambda=q/(\rho\sigma_z^2)$, so that $\vc b\herm\mt H^{-1}\vc b=\rho\sigma_z^2$;
the measured regression and the cross-wave prior are kept and the
Schur-complement condition guarantees $\mt\Sigma\succeq0$. This is the current default coherence-cap policy. Absent evidence or at zero confidence, the excitation
is diagonal. With isotropic horizontal amplitude this construction
rotates consistently under any change of the north/east heading gauge.

The mapping preserves the measured regression in the \emph{excitation
parameterization}, subject to shrinkage and the positivity cap. It need
not reproduce an arbitrary measured output cross-spectrum exactly:
\cref{eq:crossspectrum} still imposes the shared temporal shape.
Likewise, positivity alone does not ensure that the selected coherence is
physically correct for a particular hull.

\section{Coupling Confidence and Withdrawal}
\label{sec:confidence}
A second demodulation estimate with a 60~s horizon monitors the slower
prior. For each horizontal component form
$\beta_i=\hat S_{iz}/\hat S_{zz}$ with approximate regression variance
\begin{equation}
 v_i=\frac{(1-\gamma_{iz}^2)\hat S_{ii}}{N_{\rm eff}\hat S_{zz}},\qquad
 \gamma_{iz}^2=\frac{|\hat S_{iz}|^2}{\hat S_{ii}\hat S_{zz}}.
 \label{eq:regvar}
\end{equation}
Prior and monitor are compared using
\begin{equation}
 D_\beta^2=\sum_{i=x,y}
 \frac{|\beta_{i,f}-\beta_{i,p}|^2}
 {v_{i,f}+v_{i,p}+0.05^2\max(|\beta_{i,f}|^2,|\beta_{i,p}|^2)+10^{-12}}.
 \label{eq:mismatch}
\end{equation}
The overlapping estimates are correlated; these thresholds are engineering
confidence rules, not a calibrated independent chi-square hypothesis test.
Consistency is withdrawn above 16 and restored below 6.

\begin{table}[!t]
\caption{Additional default gates on shared-mode coupling. These affect
cross-axis excitation; they do not disable the marginal CARMA motion prior.}
\label{tab:gates}\centering\footnotesize
\begin{tabular}{@{}p{.28\columnwidth}p{.66\columnwidth}@{}}\toprule
Evidence & Rule\\\midrule
Warmup & Prior valid after 120~s; short monitor has 60~s history.\\
Band energy & Monitor vertical power divided by its white-noise envelope prediction: zero weight at SNR $\le10$, full at $\ge100$, log interpolation between.\\
Turning & World yaw rate smoothed over 2~s exceeds $3^\circ$/s; withdraw and hold for 30~s.\\
Period change & $|\dd\log f_{\rm op}/\dd t|$ smoothed over 10~s exceeds 0.01/s; withdraw and hold for 30~s.\\
Innovation change & 20~s mean accelerometer NIS exceeds twice its 600~s mean; recover below a ratio of 1.5.\\\bottomrule
\end{tabular}
\end{table}

The envelope noise-power approximation for the SNR gate is
$s_{n,c}^2 h\pi(0.15f_{\rm op})/2$, with separately configured coupling
noise standard deviation $s_{n,c}=0.015$~m/s$^2$. This short-monitor SNR
and the vertical amplitude evidence in \cref{eq:evidence} are different
statistics. The NIS gate uses a ratio because measurement covariance can
be deliberately inflated; an absolute expectation of three is then not
a useful reference for this monitor.

Let $e_c$ be the energy weight when all gates pass, and zero otherwise.
Confidence is smoothed every sample:
\begin{equation}
 c^+=c+(1-e^{-h/T_c})(e_c-c),\qquad
 T_c=\begin{cases}3~\mathrm s,&e_c<c,\\60~\mathrm s,&e_c\ge c.\end{cases}
 \label{eq:confidence}
\end{equation}
The applied gain changes at the approximately 0.1~s confidence cadence
and the excitation is rebuilt at the next step, with or without a model
commit. This
withdraws dependence quickly and restores it slowly. The process tends
toward diagonal excitation smoothly; ordinary smoothing does not snap
$c$ to exactly zero in finite time.

Reducing future process correlation leaves stale posterior motion
cross-covariance. Partition the remaining attitude and bias states as
$r$, horizontal motion as $h$, and vertical motion as $z$. On a confidence
decrease of at least 0.1\% since the last shrink (smaller declines
accumulate) let $\alpha_c=c_{\rm new}/c_{\rm old}$ and use
\begin{align}
 \mt E_{hz}&=\mt\Pi_{hr}\mt\Pi_{rr}^{-1}\mt\Pi_{rz},\nonumber\\
 \mt\Pi_{hz}^+&=\mt E_{hz}+
                \alpha_c(\mt\Pi_{hz}-\mt E_{hz}).
 \label{eq:shrink}
\end{align}
This removes only correlation unexplained by shared attitude and bias
uncertainty. In the conditional motion covariance it is a convex
combination of the current covariance and its horizontal/vertical block
diagonal, so it preserves positivity when the conditioning block is
positive definite. If that factorization fails, the shrink is skipped.
No estimated position, velocity or attitude is changed by this operation.


The conditional shrink has the same physical meaning in invariant coordinates:
conditioning on attitude and bias removes the common attitude shear. The
implementation uses the shared indexed Schur algorithm with the TFG
column-major indices; it does not erase all world-column cross-covariance.

\section{Parameter Commits and Covariance Transport}
\label{sec:commits}
At a model commit the four physical/spectral coordinates $[p,v,a,\ell]$ are
held. The two colouring coordinates have scale $s_k=g_0\sigma_k$ on axis $k$,
so the physical axis map is
\begin{equation}
 M_k=\diag(1,1,1,1,s_k^+/s_k^-,s_k^+/s_k^-).
\end{equation}
The nominal world coordinates satisfy
$\hat w'_{i,k}=\sum_jM_k(i,j)\hat w_{j,k}$. If these maps differ between
axes, simply scaling the same rows of invariant covariance is incorrect.
The complete invariant linearization contains an attitude shear:
\begin{equation}
 \begin{split}
 \rho'_{i,k}={}&\sum_jM_k(i,j)\rho_{j,k}\\
 &+\left((\sk{\hat w'_i})_{k,:}
       -\sum_jM_k(i,j)(\sk{\hat w_j})_{k,:}\right)\phi.
 \end{split}
 \label{eq:shear}
\end{equation}
Using \cref{eq:physicalmap}, the implementation transports covariance to
physical additive coordinates, applies the axis maps to means and the full
covariance, and returns to invariant coordinates at the new nominal state.
If $M_e$ is the corresponding full additive map, this is the congruence
\begin{equation}
 \Pi^+=L\Pi^-L\tr,\qquad L=T_{\hat X^+}M_eT_{\hat X^-}^{-1}.
\end{equation}
Every cross-axis, attitude and bias block is carried. Horizontal variance
inflation by the coherence cap uses the same operation, so it also retains
\cref{eq:shear}. A finite-difference test with nonzero colouring means and
unequal axis scales checks this map independently.

Input-normal colouring coordinates keep their unit stationary Gramian
independent of frequency scale. With compatible initialization and diagonal
excitation, the lossless identities preserve the scalar marginal CARMA law
under prescribed normalized time/amplitude scheduling at the linear-model
covariance level. They do not establish statistical consistency of the
closed-loop adaptive filter or an instantaneous stationary law after arbitrary
coupling changes. Ordinary commits never reset a posterior to a fresh stationary
prior; only explicit startup/reseed operations have that role.
\section{Startup and Transition to Live Estimation}
\label{sec:startup}
Starting a joint motion/attitude/bias filter from one acceleration sample
would identify gravity plus wave acceleration as gravity. The frontend
therefore starts first. The main filter is constructed, but the normal
wrapper holds its prediction and correction until handoff. Meanwhile the
proxy integrates gyro and acceleration, the period and energy channels
accumulate data, and magnetic acquisition seeks a reference.

\subsection{Learning a magnetic frame}
Acquisition begins no earlier than 7~s. A gravity-direction quality gate
averages world-frame specific force over 12~s, with a 5~s warmup. It
requires horizontal gravity-direction discrepancy below a sine threshold
of 0.075 for a 2~s hold, and vetoes extreme gyro motion. A separate 30~s
tilt-acquisition fallback can admit a usable reference when that gate
does not settle. This is a startup policy, not a proof of correct tilt.

Magnetometer samples are accumulated in the proxy's \emph{yaw-removed}
tilt frame. An arbitrary drifting proxy yaw must not enter the reference
being learned. At least 128 samples and a 15~s window are required by the
default acquisition configuration. The estimated reference is gauge-fixed
as $\vc B_w=[B_H,0,B_D]\tr$; a one-time yaw rotation puts the learned
magnetic north on $+X$. Its magnitude and inclination are learned, not
assumed known in advance. This north is magnetic north unless declination
is supplied elsewhere.

An initial tilt error can be embedded in that first reference. A second
acquisition starts at 90~s and averages for 30~s in the same independent
proxy tilt frame. It refines the reference after the proxy has had time
to settle. Configured hard-iron estimation is separate from the basic
wave-state model and must satisfy its own excitation and quality checks.

\subsection{Handoff and bias release}
The normal handoff requires an initialized proxy, at least 8~s elapsed,
held gravity agreement, acquired north when magnetic aiding is enabled,
and a ready tuner with a usable measured period. The latter ordinarily
prevents normal handoff before roughly 31.8~s. A 150~s timeout uses the
best available attitude but still requires the aligned gravity branch;
it does not authorize the antipodal tilt. Magnetic acquisition timing can
extend the applicable timeout floor. Thus timeout entry is a degraded
startup path, not evidence that every quality condition was satisfied.

The nominal attitude is seeded from the proxy tilt composed with the
learned yaw gauge. Handoff standard deviations are 0.035~rad for tilt and
0.087~rad for magnetically gauged yaw; without a yaw gauge the latter is
1.5708~rad. The current process scale is applied, and the motion covariance
is seeded from its stationary prior with unrelated cross-blocks cleared.
This initialization is different from the later parameter commits, which
carry the posterior forward.

Accelerometer-bias adaptation remains held at handoff. The default unlock
requires 250 actually accepted magnetic corrections spanning more than
one second, and Live status. A correction counted as accepted by the core
still does not certify strong geometric information or an undisturbed
field. Holding bias initially reduces the chance that it absorbs a
settling tilt error. It does not remove the physical need for subsequent
magnetic service.

The coupling prior and short monitor receive data only once main-filter
accelerometer updates run. Their 120~s and 60~s clocks therefore begin
after handoff in the normal wrapper, not at power-on. With default
confidence initially zero, initial Live estimation has diagonal excitation;
shared coupling is acquired gradually later. An optional presentation
detrender can process output displacement, but is not an $S=0$ update
and is not part of the 27-state estimation model.

\section{Per-Sample Execution}
\label{sec:execution}
The order of operations is part of the adaptive algorithm. During
bootstrap only the measurement frontend and magnetic acquisition run.
In Live operation one IMU sample is processed as follows:
\begin{enumerate}
 \item Apply any staged tune from the preceding sample; refresh motion
 dynamics, paired covariance integrals and colouring-coordinate transport
 when required.
 \item Condition the measured specific force, including a configured
 lever arm, and advance the private attitude proxy. Update the configured
 vibration/low-wave measurement-noise conditioning.
 \item Predict the main attitude, motion and biases over the actual
 sample interval. Propagate the complete covariance.
 \item Reconstruct measured world acceleration for both cross-spectral
 accumulators. Perform the accelerometer update and full group reset.
 Update coupling confidence using the new innovation diagnostic.
 \item Run the tilt watchdog; advance the vertical acceleration tracker,
 energy statistics and tuner. Stage a tune when the commit clock is due.
 Advance the independent wave-period estimator for subsequent samples.
\end{enumerate}
Magnetic samples arrive on their own clock. They acquire or refine the
reference during startup and deliver vector corrections in Live
operation. Coupling confidence withdrawal can shrink conditional motion
covariance at its own approximately 0.1~s cadence, while the new excitation
is realized on the next motion-model refresh. The acceleration-band
frequency tracker also serves wave-direction reporting; its carrier is
not substituted for the zero-crossing period in \cref{eq:kappa}.


In the TFG core, startup covariance seeding, bias holds and explicit attitude
writes are defined in physical coordinates and mapped through
\cref{eq:physicalmap}. This preserves their intended physical uncertainty
rather than assigning a numerically identical but physically different
matrix in the invariant tangent frame.

\input{underway-baseline.tex-part}

\section{Deterministic Replay and Quality Gates}
\label{sec:validation}
The dedicated simulator replays the pinned
\texttt{oceanography-waves-lib v1.3.0} 5 kn underway vessel-response records,
with heading oscillating $\pm3^\circ$ about north every 120 s. There are eight
scored records: JONSWAP and PM--Stokes, each at incident significant wave
heights 0.27, 1.5, 4.0 and 8.5~m. Each history lasts 1200~s at 200~Hz. The
sensor and initialization seeds, simulated noise, magnetic schedule and
frontend are the same as the shared-mode quaternion replay. The simulated
wave realizations are fixed; these eight deterministic results are neither
independent sea trials nor population confidence intervals.

The first 300~s are excluded from the scored window. Statistics use all
180,000 remaining samples; rendering decimation never changes the gate or
reported RMS. Displacement X/Y/Z are normalized by incident $H_s$. The 3-D
error is normalized by the maximum norm of the reference displacement in the
same trailing window. The accelerometer Z-bias percentage uses the maximum
absolute true Z-bias component; the 3-D accelerometer and gyroscope bias
percentages use the maximum norm of the corresponding true residual
bias vector, all over that window. These denominators are different and must
not be interpreted as interchangeable percentages.

\input{tfg-shared-sim-results-generated.tex-part}
\input{tfg-shared-gates-generated.tex-part}

The executable's limits are its own measured worst errors over these eight
records, enlarged by 0.5\% and rounded upward to 0.00001. The limits are locked
regression sentinels, not estimated uncertainty bounds and not tuned anew at
each CI run. The baseline JSON pins input hashes, compiler, source revision,
per-record measurements and the baseline executable hash. A regression change
must be assessed against that fixed baseline rather than silently relaxing
its limits.

The test script verifies the pinned input hashes, requires eight actual
passing quality gates, and rejects empty or truncated replays. Nonfinite
motion, attitude or bias errors fail instead of bypassing comparisons.
The explicit regression mode rejects protocol overrides that would invalidate
the limits. A negative test corrupts only displacement truth while retaining
the delivered inertial samples; its large RMS error must trip the actual
numeric gates and return a failing status. This tests gate activation as well
as the ordinary passing path. Alternative seeds, magnetometer outages,
noise-free runs and coupling ablations remain useful diagnostics but do not
establish passage of this fitted default-seed regression.

\begin{figure}[!t]\centering
\IncludePGF{w3d_tfg_csm_method_jonswap_medium_xykin.pgf}
\caption{Surge and sway tracking of the current deterministic medium JONSWAP
record. Drawing shows the final 60~s; annotations use the full final 900~s.}
\end{figure}
\begin{figure}[!t]\centering
\IncludePGF{w3d_tfg_csm_method_jonswap_medium_zkin.pgf}
\caption{Heave, vertical velocity and acceleration of the same gated replay.}
\end{figure}
\begin{figure}[!t]\centering
\IncludePGF{w3d_tfg_csm_method_jonswap_medium_att.pgf}
\caption{Roll, pitch and heading from the same deterministic replay.}
\end{figure}
\begin{figure}[!t]\centering
\IncludePGF{w3d_tfg_csm_method_jonswap_medium_gyro_bias.pgf}
\caption{Gyroscope bias tracking across the full replay, including startup.}
\end{figure}
\begin{figure}[!t]\centering
\IncludePGF{w3d_tfg_csm_method_jonswap_medium_acc_bias.pgf}
\caption{Accelerometer bias tracking. Position accuracy alone does not establish
unique identification of sensor bias and motion.}
\end{figure}

\section{Comparisons, Vibration and Sea-State Changes}
\label{sec:scope}
The comparative companion uses a paired $2\times2$ design: quaternion and
TFG error geometry, each with independent-axis and shared-mode CARMA priors.
Its primary synthetic directional-sea ensemble has 24 sensor/initialization
seeds in 27 conditions, 1200~s at 200~Hz, scored from 300~s. Every filter in
a pair receives the same physical history and corruption. The figures in
this section reuse that committed paired-study evidence and its compact
traces; the ensemble has not been rerun merely to write this method article.
The deterministic table above is separately regenerated by the current replay.

\begin{figure*}[!t]\centering
\IncludePGF{../../reports/results/tfg_shared_modes/figures/paired_position.pgf}
\caption{Paired position changes of shared TFG relative to shared quaternion
estimation. Negative values favor TFG; intervals use independent seeds.
Geometry and stochastic-prior benefits are distinct effects.}
\end{figure*}
The companion's condition-balanced primary comparisons favor shared TFG in
position RMS, but the geometric increment is small compared with changing
from independent to shared excitation. This does not make the method uniformly
best in every condition or resolve the poor joint error calibration reported
by the companion. Bias/motion ambiguity and process mismatch survive a change
of error coordinates. Reduced-heave response is a particularly important
limitation of a directional shared prior; the coherence-cap policy limits its
statistical dependence instead of assuming a universal hull response.

\begin{figure}[!t]\centering
\IncludePGF{../../reports/results/tfg_shared_modes/figures/engine_degradation.pgf}
\caption{Machinery-vibration degradation from the paired synthetic study.
Error bars reflect seed variation. The engine-off case uses paired seeds.}
\end{figure}
Engine vibration adds accelerometer and gyro corruption outside the desired
wave motion. A conditioning filter and measurement covariance inflation can
reduce its influence, but neither guarantees separation of wave energy from
installation vibration. The frequency content, sensor bandwidth, mounting and
rectification matter; the plotted engine model is a declared simulation,
not a calibration for every diesel or boat.

\begin{figure*}[!t]\centering
\IncludePGF{../../reports/results/tfg_shared_modes/figures/roundtrip.pgf}
\caption{Sea-state growth and decay in the paired study. The physical amplitude
scale moves from 0.5 to 4 at 400~s and back at 800~s through differentiated
20~s smooth transitions. No truth scale is supplied to the operating tuner.}
\end{figure*}
Changing sea state tests both tightening and release of the adaptive prior.
Period and energy statistics respond at their own horizons; staged commits
carry the posterior rather than reseeding it. Fast confidence withdrawal
reduces stale cross-axis dependence during a mismatch, while slower recovery
requires renewed supporting evidence. Full-rate transient metrics precede
the peak-preserving drawing reduction. The companion additionally examines
abrupt direction changes and magnetic outage/reacquisition; these tests do
not imply continued absolute yaw observability during an outage.

\section{Computational and Practical Limits}
The covariance contains 729 scalar entries before symmetry is exploited.
The six world columns permit structured propagation, with the gyro-noise
correlations and full reset retained. Fixed-size member scratch avoids a
covariance-sized real-time stack temporary. The shared implementation avoids
a second spectral fit, a second adaptation path or a copied confidence monitor.
It still costs more than the smaller independent-axis TFG and than shared
quaternion estimation; reduced state count relative to 33 does not establish
an embedded 200~Hz deadline. The companion's Xtensa sizes and compiler stack
bounds are build-time evidence, not measured stack high-water marks or
on-device timing.

The local-equilibrium displacement model excludes permanent position offsets
and current. A single shared temporal shape cannot represent every crossing
sea, hull transfer function, manoeuvre or low-frequency acceleration. Spectral
fit quality is band limited. Gaussian tangent covariance and finite group
corrections are approximations; process-law equivalence at zero coupling does
not imply bitwise equality of nonlinear filters with different latent spaces.
Magnetic reference error and unqualified disturbance are still important failure
modes. Sea trials, device-specific noise measurements, and installation tests
remain necessary before treating simulation accuracy as operational performance.

\section{Conclusion}
The method combines a stable CARMA displacement prior, complex directional
excitation and right-invariant two-frame estimation. Minimal realization gives
18 real motion coordinates and nine attitude/bias errors. The prior regularizes
wave-relative displacement without an integral pseudo-measurement; the independent
frontend measures its operating scale, and magnetic aiding supplies complementary
orientation information. A full group reset and physical-coordinate transport
carry the covariance through finite corrections and anisotropic adaptive changes.
The dedicated replay and explicit gates make changes to this implemented method
reviewable on pinned vessel-response data. Those checks complement, rather than
replace, the comparative study and deployment validation.

\FloatBarrier\appendices
\section{Constructing the Two-State Colouring Filter}
\label{app:colour}
At unit scale write $P_+(s)=s^2+b_1s+b_0$ with complex coefficients.
Then $N_r=s^2+\Rea b_1s+\Rea b_0$ and
$N_i=\Ima b_1s+\Ima b_0$. The even polynomial on the right of
\cref{eq:D} is $\omega^4+\beta\omega^2+\gamma_0$, where
\begin{align}
 \gamma_0&=|b_0|^2,\quad \beta=|b_1|^2-2\Rea b_0,\nonumber\\
 d_0&=\sqrt{\gamma_0},\quad d_1=\sqrt{\beta+2d_0}.
 \label{eq:Dcoeff}
\end{align}
Positive $d_0,d_1$ select the stable real spectral factor. An explicit
observer-form realization of \cref{eq:V} is
\begin{align}
 \mt A_0&=\begin{bmatrix}-d_1&1\\-d_0&0\end{bmatrix},\quad
 \mt C_0=[1\quad0],\nonumber\\
 \mt B_0&=\begin{bmatrix}\Rea b_1-d_1&\Ima b_1\\
                         \Rea b_0-d_0&\Ima b_0\end{bmatrix}.
 \label{eq:observercolour}
\end{align}
Solve $\mt A_0\mt X+\mt X\mt A_0\tr+\mt B_0\mt B_0\tr=0$,
factor $\mt X=\mt L\mt L\tr$, and transform with $\mt L^{-1}$.
An orthogonal two-state rotation then makes $\mt B_w$ upper triangular
with positive diagonal, fixing otherwise arbitrary coordinate freedom.
The result satisfies \cref{eq:inputnormal}. The implementation uses
closed-form two-by-two algebra for this construction, once per shape,
then uses \cref{eq:colourscale} at each frequency scale. A generic
factorization in every sample is unnecessary.

\section{Noise Moments Used by the Evidence Channel}
\label{app:evidence}
Put $c_n=s_n^2h$, $q'=q_b/(ab)$ and let $w$ be the accumulated normalized
EMA weight. The implemented effective degrees of freedom are
\begin{equation}
 \nu=2T_V(a+b)\frac{(c_n+q')^2}
 {c_n^2+2c_nq'+q'^2(a^2+3ab+b^2)/(ab)}\frac{w}{2-w}.
 \label{eq:nu}
\end{equation}
The weight correction accounts for the shorter record immediately after
startup. Under white noise alone the long-record expression reduces to
$2T_V(a+b)$. The extra low-pass and derivative in \cref{eq:fc} use
the continuous-time no-wave moment approximations
\begin{align}
 V_{y_2,0}&=\frac{c_n b^3}{4(a+b)^2}
       +\frac{q_b b(a+2b)}{4a(a+b)^2}-\frac{q_b}{2T_Va^2},\nonumber\\
 V_{\dot y_2,0}&=\frac{c_n b^4(2a+b)}{4(a+b)^2}
                       +\frac{q_b b^3}{4(a+b)^2}.
 \label{eq:noisemoments}
\end{align}
Finite discrete differences, changing band corners and non-white
installation noise limit these approximations. The uncertainty term $u$
in \cref{eq:evidence} prevents treating the nominal noise prediction as
an exact known quantity.


\section{Implementation and Reproduction}
The core is \path{src/kalman_carma_shared_modes/KalmanTfgCarmaSharedModes.h};
its common two-frame engine is
\path{src/kalman_carma/Kalman3D_Wave_TFG_CARMA3.h}. The minimal motion and
confidence components are shared with quaternion estimation. The public wrapper
is \path{SeaStateFusion_TfgCarmaSharedModes}; it uses the common CARMA schedule,
wave evidence and magnetic startup. The group equations here use body-to-world
$R$; the public quaternion accessors retain the project's world-to-body convention.

The dedicated replay is
\path{tests/kalman_tfg_shared_modes/kalman_tfg_shared_modes-sim.cpp}.
Both geometries reuse \path{tests/common/SharedModesReplay.h}; TFG limits are in
\path{reports/results/filter_sim_gates/limits.json}.
Independent mathematical checks are in
\path{tests/kalman_carma_shared_modes/tfg_shared_modes-test.cpp}.

From the repository root, run \texttt{make all}, or
\texttt{make -C tests/kalman\_tfg\_shared\_modes all} for this replay.
Regenerate tables with \path{tools/tfg_shared_modes_sim_table.py} and figures
with \path{plots/kalman_tfg_shared_modes/draw_plots.sh}. Stage PGFs as CI does:
\begin{verbatim}
cp plots/kalman_tfg_shared_modes/*.pgf \
   doc/kalman_tfg_shared_modes/
\end{verbatim}
Build with \texttt{latexmk -lualatex} on
\path{doc/kalman_tfg_shared_modes/kalman_tfg_shared_modes-method.tex} from its
directory.

The comparative study remains \path{kalman_tfg_shared_modes.tex}; the complete
per-record gallery is \path{kalman_tfg_shared_modes-charts.tex}. The fit and
paired experiments use \path{tools/carma_fit.py} and
\path{tools/tfg_shared_modes_study.py}. Statistics use full-rate data;
drawing uses bounded, peak-preserving decimation.
Source and declared evidence are available at
\url{https://github.com/bareboat-necessities/ocean-imu}.

\section*{Acknowledgments}
ChatGPT assisted with source inspection, mathematical drafting and editorial
preparation under author supervision. The author is responsible for the method,
validation and final manuscript.
\begin{thebibliography}{9}
\bibitem{jonswap}
K. Hasselmann \emph{et al.}, ``Measurements of wind-wave growth and swell
decay during the Joint North Sea Wave Project (JONSWAP),''
\emph{Erg\"anzungsheft zur Deutschen Hydrographischen Zeitschrift},
Reihe A (8), no. 12, 1973.
\bibitem{fasen}
V. Fasen-Hartmann and M. Scholz, ``Factorization and discrete-time
representation of multivariate CARMA processes,'' arXiv:2102.11681, 2021.
\bibitem{vanloan}
C. F. Van Loan, ``Computing integrals involving the matrix exponential,''
\emph{IEEE Trans. Autom. Control}, vol. 23, no. 3, pp. 395--404, 1978.
\bibitem{mahony}
R. Mahony, T. Hamel and J.-M. Pflimlin, ``Nonlinear complementary filters
on the special orthogonal group,'' \emph{IEEE Trans. Autom. Control},
vol. 53, no. 5, pp. 1203--1218, 2008.
\bibitem{barrau}
A. Barrau and S. Bonnabel, ``The geometry of navigation problems,''
\emph{IEEE Transactions on Automatic Control}, 2022,
doi:10.1109/TAC.2022.3144328.
\end{thebibliography}
\end{document}
